\section{Introduction}
\label{sec:intro}

Auto-labelling made large segmentation benchmarks possible. It also made some of
them circular, because the label is a deterministic function of the image the model
receives. One measurable form of that circularity occurs when the labeller is run on
crops and its answer changes with the crop. That failure is widely acknowledged and
rarely quantified, in part because there has been no estimand with a value known in
advance under the null that can be computed on a benchmark you did not build.

This paper supplies one. We define crop alignment, prove its null, calibrate it
against a labeller we control, measure what the thing it detects costs against human
annotation, and release an implementation that runs on other people's data. We apply
it to two benchmarks, and \Cref{sec:wild,sec:resolution} report what it finds there.

Foundation-model masks are increasingly used as ground
truth~\cite{kirillov2023segment}. When those masks are generated crop by crop, they
can inherit the property tested here; scene-level pseudo-labels can still close the
input--label loop but fall outside our diagnostic if they are crop-invariant. The
adjacent literatures in \Cref{sec:related} diagnose related failures after the fact.
What is scarce is a controlled measurement of this crop-dependent case.

\paragraph{Terms.} A benchmark is \emph{closed-loop} when its labels are computed
from the model's own input. \emph{Crop alignment} ($\kap$) is defined in
\Cref{sec:kappa} and audits only the crop-dependent subset of such benchmarks. The
\emph{artefact premium} is defined in
\Cref{sec:resolution}, where it is also estimated.
